\chapter*{Resumen}
\addcontentsline{toc}{chapter}{Resumen}

Los ecosistemas de criptoactivos y finanzas descentralizadas integran exchanges centralizados (CEX), exchanges descentralizados (DEX), Automated Market Makers (AMM), protocolos de préstamo, oráculos, liquidadores, proveedores de liquidez y agentes automatizados. Estas interacciones combinan reglas de protocolo, infraestructura distribuida, costos de ejecución e incentivos económicos, por lo que se comportan como sistemas complejos adaptativos. Sin embargo, una parte importante de la literatura y de las herramientas disponibles se concentra en predicción de precios o backtesting simplificado y deja en segundo plano la trazabilidad de datos, la microestructura híbrida, la liquidez fragmentada, el slippage, el gas, la latencia, las liquidaciones y los riesgos de integridad de mercado.

Esta tesis propone el diseño de una infraestructura experimental multi-agente para modelar, simular y analizar escenarios CEX/DEX de forma reproducible. La infraestructura organiza cinco capas: adquisición y validación de datos, representación de variables de mercado y protocolo, simulación de agentes heterogéneos, escenarios de estrés y evaluación mediante artefactos versionados. El propósito no es producir recomendaciones de inversión ni un bot rentable, sino habilitar experimentos auditables sobre liquidez, ejecución, arbitraje con latencia, cascadas de liquidación y propagación de shocks.

La versión \versionTesis{} presenta el problema, el marco conceptual, el diseño metodológico y el protocolo de evaluación. Los resultados empíricos se incorporarán en versiones posteriores únicamente después de contar con datasets documentados, código versionado y ejecuciones reproducibles.

\textbf{Palabras clave:} sistemas complejos, simulación basada en agentes, CEX, DEX, DeFi, microestructura de mercados, reproducibilidad, blockchain.

\clearpage
\selectlanguage{english}
\chapter*{Abstract}
\addcontentsline{toc}{chapter}{Abstract}

Crypto-asset and decentralized-finance ecosystems connect centralized exchanges (CEX), decentralized exchanges (DEX), Automated Market Makers (AMM), lending protocols, oracles, liquidators, liquidity providers, and automated agents. Their interactions combine protocol rules, distributed infrastructure, execution costs, and economic incentives; therefore, they can be understood as complex adaptive systems. Yet much of the available literature and tooling focuses on price prediction or simplified backtesting, while data traceability, hybrid market microstructure, fragmented liquidity, slippage, gas, latency, liquidations, and market-integrity risks remain secondary.

This thesis proposes the design of a reproducible multi-agent experimental infrastructure to model, simulate, and analyze CEX/DEX scenarios. The infrastructure organizes five layers: data acquisition and validation, market and protocol variables, heterogeneous-agent simulation, stress scenarios, and evaluation through versioned artifacts. Its purpose is not to produce investment advice or a profitable trading bot, but to enable auditable experiments on liquidity, latency-aware arbitrage, liquidation cascades, execution, and shock propagation.

Version \versionTesis{} presents the problem, conceptual framework, methodology, and evaluation protocol. Empirical results will be added only in later versions after documented datasets, versioned code, and reproducible executions are available.

\textbf{Keywords:} complex systems, agent-based simulation, CEX, DEX, DeFi, market microstructure, reproducibility, blockchain.
\selectlanguage{spanish}
\clearpage
